% ============================================================
% ch08_09.tex —— 第 8 章、第 9 章
% 底本：Springer LNM 2386, printed pp.77–92 = PDF pp.117–132
% 公式唯一真值：work/pages/pNNN.png（200 dpi）
% ============================================================

% [原文 p.77]
\chapter{黎曼曲面上的微分}

\ChapterAbstract{我们回顾亚纯微分的基本事实，这些事实在后续谱网络的解析构造中将会用到。特别地，二次微分诱导出曲面的自然分解，这些分解给出三角剖分，从而有助于理解曲面的某些微分几何面向。关于二次微分在泰希米勒理论中作用的标准参考文献包括 [162] 或 [336]。}

\section{阿贝尔微分及其轨道}

设 $X$ 是紧黎曼曲面，其底拓扑曲面为 $S$。由于 $X$ 的复维数按定义为 1，其上定义有次数至多 2 的微分形式。微分形式在坐标变换下不变，因而在整个空间 $X$ 上有定义。于是我们可以给出如下定义：

\begin{definition}[8.1] 设 $X$ 是黎曼曲面，并取 $X$ 上的局部坐标 $z$。$X$ 上的阿贝尔微分局部地写为 1-形式 $\omega=f(z)\dd z$，其中 $f$ 是 $X$ 上的光滑函数。换言之，它是 $K$ 的光滑截面，其中 $K=\Omega^{1,0}(S)$ 表示典范丛。若局部定义微分 $\omega$ 的每个函数都是全纯的（相应地，亚纯的），则称 $\omega$ 为全纯阿贝尔微分（相应地，亚纯阿贝尔微分）。\end{definition}

下文除另有说明外，我们考虑亚纯阿贝尔微分。

Riemann–Roch 定理的一个应用表明，$X$ 上全纯阿贝尔微分构成的向量空间 $\mathrm{H}^0(X,\Omega)$ 的维数等于黎曼曲面的亏格 $g$。这可用解析方法证明（例如见 [126]）；对某些特定情形的黎曼曲面，[246] 中给出了全纯阿贝尔微分的一组基 $\omega_1,\dots,\omega_g$ 的显式描述。

$f(z)$ 的零点与极点不依赖于局部坐标的选取，因此可以分别定义相应亚纯微分的零点与极点。

% [Tutor Note: 「$\dim\mathrm{H}^0(X,\Omega)=g$」的核对方法：对 $D=K$ 用 Riemann–Roch 得 $h^0(K)-h^0(K^{-1})=\deg K-g+1=(2g-2)-g+1=g-1$；又 $h^0(K^{-1})=h^0(\mathcal{O}_X)=1$（常数函数），故 $h^0(K)=g$。易错点：不能漏掉 $h^0(\mathcal{O}_X)=1$ 这一项，否则会误得 $g-1$。]

% [原文 p.78]
对奇异点 $z_0\in X$ 附近的 Laurent 展开
\[
f(z)=\sum_{n=-\infty}^{\infty}\alpha_n(z-z_0)^n
\]
而言，阿贝尔微分 $\omega$ 在点 $z_0$ 处的留数 $\operatorname{Res}_{z_0}\omega$ 定义为该展开中的系数 $\alpha_{-1}$。留数同样不依赖于局部坐标的选取，因为它可以写为
\[
\operatorname{Res}_{z_0}\omega=\frac{1}{2\pi\ii}\int_{\partial B_{z_0}}\omega(z)\,,
\]
其中 $B_{z_0}$ 是点 $z_0$ 周围的圆盘，使得微分形式 $\omega$ 在 $B_{z_0}$ 的闭包内没有别的奇异点。注意，对 $X$ 上微分 $\omega$ 的有限个奇异点 $z_0,\dots,z_n$，有 $\sum_{j=0}^{n}\operatorname{Res}_{z_j}\omega=0$。事实上，对每个奇点 $z_j$ 取闭包内不含其他奇点的圆盘 $B_j$，即见
\[
\sum_{j=0}^{n}\operatorname{Res}_{z_j}\omega=\frac{1}{2\pi\ii}\sum_{j=0}^{n}\int_{\partial B_{z_j}}\omega(z)=-\frac{1}{2\pi\ii}\int_{X\setminus\cup B_{z_j}}\dd\omega(z)=0\,,
\]
因为微分 $\omega$ 在 $X\setminus\cup B_{z_j}$ 上全纯，故在其上有 $\dd\omega=0$。

% [Tutor Note: 上式中间的负号来自边界定向。以 $X\setminus\cup_j B_{z_j}$ 为区域用 Stokes 公式时，其边界是 $-\cup_j\partial B_{z_j}$（内边界取反向），故 $\sum_j\int_{\partial B_{z_j}}\omega=-\int_{X\setminus\cup_j B_{z_j}}\dd\omega$。可核对的例子：在 $\C\PP^1$ 上取 $\omega=\dd z/z$，$z=0$ 与 $z=\infty$ 处留数分别为 $+1$ 与 $-1$，和为 $0$。]

\begin{definition}[8.2] 设 $X$ 是黎曼曲面，$\omega$ 是 $X$ 上的阿贝尔微分。$\omega$ 的临界点是指它的零点与极点，$X$ 上所有其他点称为正则点。注意对任何阿贝尔微分 $\omega$，临界点只有有限多个。\end{definition}

现设 $\omega$ 是一般非零的阿贝尔微分，以 $P$ 记 $\omega$ 的极点集，以 $Z$ 记 $\omega$ 的零点集。$\omega$ 的虚部 $\operatorname{Im}(\omega)$ 是 $X$ 上的实 1-形式场。由于 $X$ 的切丛实秩为 2，形式 $\operatorname{Im}(\omega)$ 一般地对每个 $z\in X$ 有 1 维核 $\operatorname{Ker}(\operatorname{Im}(\omega_z))\subset T_zX$。这样的对象通常称为一个分布。更确切地说，$X$ 的切丛 $TX$ 中的 1 维分布是 $TX$ 的 1 维子丛。在我们的情形，形式 $\omega$ 只能有孤立零点，故分布 $\operatorname{Ker}(\operatorname{Im}(\omega))$ 在 $X':=X\setminus(P\cup Z)$ 上有良好定义。

分布 $\operatorname{Ker}(\operatorname{Im}(\omega))$ 在 $X'$ 上给出一个叶状结构 $\hat{\mathcal{F}}_\omega$，即一族光滑无向曲线（称为该叶状结构的叶），它们在每一点都与 $\operatorname{Ker}(\operatorname{Im}(\omega))$ 相切，并且 $X'$ 的每一点恰有 $\hat{\mathcal{F}}_\omega$ 的一片叶。我们可以给 $\hat{\mathcal{F}}_\omega$ 指定定向：在每点 $z\in X\setminus Z$ 光滑地选取方向 $\R_+v_z\subset\operatorname{Ker}(\operatorname{Im}(\omega))$，使得 $\operatorname{Re}(\omega(v_z))>0$。定向后的叶状结构记作 $\mathcal{F}_\omega$（图 8.1）。

注意，$\mathcal{F}_{-\omega}$ 可由 $\mathcal{F}_\omega$ 反转定向得到。

\begin{example}[8.3] 取 $X=\C$，考虑两个阿贝尔微分 $\omega_1=\dd z$ 与 $\omega_2=(z-z_0)^2\dd z$，其中 $z_0\in X$。则在 $z_0$ 的小邻域内，叶状结构 $\mathcal{F}_{\omega_1}$ 与 $\mathcal{F}_{\omega_2}$ 如图 8.1 所示。\end{example}

% [原文 p.79]
\Fig[0.40\linewidth]{figures/fig_8_1.png}{8.1}{围绕 $z_0$ 的叶状结构 $\mathcal{F}_{\omega_1}$（上）与 $\mathcal{F}_{\omega_2}$（下）。}

\section{二次微分及其轨道}

除了考虑作为 1-形式的阿贝尔微分，我们还可以考虑更高阶 $k\in\N_{\ge 2}$ 的微分。它们是局部形如 $f(z)\dd z^k$ 的对称张量，其中 $f$ 是局部光滑函数。特别有趣的是二次微分的情形，这是我们接下来要研究的。

\begin{definition}[8.4] 设 $X$ 是黎曼曲面，其底拓扑曲面为 $S$，并取 $X$ 上的局部坐标 $z$。$X$ 上的二次微分局部地写为 $\omega=f(z)\dd z^2$，其中 $f$ 是 $X$ 上的光滑函数。换言之，它是 $K\otimes K$ 的光滑截面，其中 $K=\Omega^{1,0}(S)$ 表示典范丛。若局部定义微分 $\omega$ 的每个函数都是全纯的（相应地，亚纯的），则称 $\omega$ 为全纯二次微分（相应地，亚纯二次微分）。\end{definition}

下文除另有说明外，我们考虑亚纯二次微分。

% [原文 p.80]
\begin{remark}[8.5] 与阿贝尔微分的情形不同，二次微分的轨道没有标准的定向选取。事实上，对亚纯二次微分 $\omega$ 同样可以考虑分布 $\operatorname{Ker}(\operatorname{Im}(\omega))$，它在 $X'$ 上给出叶状结构 $\hat{\mathcal{F}}_\omega$。然而不能给 $\hat{\mathcal{F}}_\omega$ 指定标准定向。这是因为对每个 $z$ 与每个 $0\ne v\in\operatorname{Ker}(\operatorname{Im}(\omega))$，$\operatorname{Re}(\omega_z)(v)$ 的符号并不依赖于 $v$，而由 $\operatorname{Re}(f(z))$ 的符号给出。下文将会看到，事实上存在某些二次微分，其相应的无向叶状结构 $\hat{\mathcal{F}}_\omega$ 不可定向。\end{remark}

Riemann–Roch 定理的一个经典应用给出：在亏格至少为 2 的闭黎曼曲面上，$X$ 上全体全纯二次微分构成的空间 $Q(X)$ 的复维数为 $3g-3$；证明可参见例如 [218, Cor. 5.4.2]。事实上，如 [370] 利用适当调和映射的 Hopf 微分所示，泰希米勒空间 $\mathcal{T}(S_g)$ 同构于向量空间 $Q(X)$。

\begin{definition}[8.6] 设 $\omega$ 是黎曼曲面 $X$ 上的二次微分，局部写为 $\omega=f(z)\dd z^2$。对 $\theta\in[0,2\pi)$，关于 $\omega$ 的直 $\theta$-弧是指光滑曲线 $\gamma$，沿其有 $\arg(\omega)=\theta\bmod 2\pi$。若 $\gamma$ 是极大的，即 $\gamma$ 不真包含于任何其他直弧之中，则称曲线 $\gamma$ 为 $\theta$-轨道。\end{definition}

注意直弧只含二次微分 $\omega$ 的正则点。我们将关注所谓的水平轨道，即 0-轨道，以及竖直轨道，即 $\frac{\pi}{2}$-轨道。下面的定理给出极大轨道的存在性与唯一性：

% [Tutor Note: 定理 8.7 的局部模型可显式写出：在正则点的小邻域内取 $w=\int\sqrt{\omega}$（沿任一路径积分），则 $\omega=\dd w^2$；关于 $\omega$ 的直 $\theta$-弧正是直线 $\arg w=\theta/2$ 的原像，故过每一点恰有一条极大直弧，存在唯一性归结为常微分方程解的存在唯一性定理。]

\begin{theorem}[8.7（{{{{{[336, Thm. 5.5]}}}}}）] 设 $\omega$ 是黎曼曲面 $X$ 上的二次微分。则过 $\omega$ 的每个正则点，存在唯一一条极大轨道。\end{theorem}

由该定理可知，极大轨道可由一个开区间参数化。于是可以考虑参数化区间两个端点处的两个极限。当这样的极限存在时，我们说它定义了轨道的一个端。注意极大轨道的端不可能是正则点，否则就与极大性矛盾。

结合注 8.5，注意人们也常把二次微分的轨道称为叶状结构 $\hat{\mathcal{F}}_\omega$ 的叶。图 8.2 展示了叶状结构的水平叶与竖直叶在单零点附近的行为，图 8.3 则类似地展示亚纯二次微分在各种阶的极点附近的行为。

\section{阿贝尔微分与二次微分的临界图}

对阿贝尔微分或二次微分 $\omega$，无向叶状结构 $\hat{\mathcal{F}}_\omega$ 的叶所携带的信息可以汇总为如下定义的临界图这一概念。

% [原文 p.81]
\Fig[0.36\linewidth]{figures/fig_8_2.png}{8.2}{单零点附近的叶状结构 $\hat{\mathcal{F}}_\omega$。}

\Fig[0.66\linewidth]{figures/fig_8_3.png}{8.3}{不同阶极点附近的叶状结构 $\hat{\mathcal{F}}_\omega$。}

% [Tutor Note: 图 8.3 图内标号（Pole of order 1、Different poles of order 2、k asymptotic lines、k=1, pole of order 3、k=3, pole of order 5）是原书图内英文，未改；图注只译原图注文字。]

\begin{definition}[8.8] 设 $\omega$ 是 $X$ 上的阿贝尔或二次亚纯微分。若 $\hat{\mathcal{F}}_\omega$ 的一片叶的某一端是 $\omega$ 的零点，则称该叶为临界的。若临界叶的另一端是 $\omega$ 的极点，则称该临界叶为正则的。若临界叶的两端都是 $\omega$ 的零点，则称之为鞍连接。此外，若一片叶在 $X$ 中的闭包具有正测度，则称它为回归的。$\omega$ 的临界图是指 $\hat{\mathcal{F}}_\omega$ 的全部临界叶的集合。若二次微分的所有临界叶都是正则的，则称之为一般的。\end{definition}

% [Tutor Note: 定义 8.8 中的「临界叶」取的是无向叶状结构 $\hat{\mathcal{F}}_\omega$ 的叶；同一几何叶在定向叶状结构 $\mathcal{F}_\omega$ 下对应两条方向相反的定向叶，故讨论临界图时不必区分定向。核对时注意三点互不重叠：正则临界叶的一端必为极点，鞍连接两端皆为（单）零点，而「回归叶」不是临界叶的补充，它按闭包的测度独立定义。]

设 $\omega$ 是一般的亚纯二次微分，至少有一个阶数不小于 2 的极点，且只有单零点。则 $\omega$ 定义一个三角剖分。更确切地说，对 $\omega$ 的每个零点，恰有三片叶终止于极点。对每两个这样的极点，都有叶状结构的叶连接它们。

% [原文 p.82]
对上述每一对极点，我们固定其中的一片叶，从而得到 $X$ 的一个三角剖分（见图 8.4）。

\Fig[0.56\linewidth]{figures/fig_8_4.png}{8.4}{由亚纯二次微分的叶状结构给出的理想三角剖分。}

Strebel 的一个基本定理 [336, Thm. 23.5] 现在给出如下结论；下文所采用的形式见 [332, Thm. 2.1]：

\begin{theorem}[8.9] 给定亏格 $g\ge 0$ 的紧黎曼曲面 $X$、$n\ge 1$ 个标记点 $p_i$ 以及 $n\ge 1$ 个正整数 $a_i$（$i=1,\dots,n$），且满足 $2g-2+n>0$，则 $X$ 上存在唯一的亚纯二次微分 $\omega$ 满足以下条件：
\begin{enumerate}[label=\arabic*.]
\item 微分 $\omega$ 在每个标记点 $p_i$ 处有二重极点，留数为 $a_i$，且没有其他极点。
\item $\omega$ 的每片闭叶都恰好绕其中一个标记点。
\item 二次微分 $\omega$ 没有回归叶。
\end{enumerate}
\end{theorem}

黎曼曲面上一般的二次微分具有水平叶与竖直叶，它们自然地给出把曲面切成欧氏矩形的方式。水平轨道与竖直轨道连同该二次微分一起确定两个相互横截的测度叶状结构。若在局部坐标 $z$ 中，存在 $C^1$-拓扑结构分别等价于 $z^r\dd z^2$（某个整数 $r\ge 1$）的水平叶与竖直叶，则称两个叶状结构在一个奇点处横截。

% [Tutor Note: 定理 8.9 的计数可作量纲核对：紧黎曼曲面 $X$ 的模空间维数为 $3g-3+n$，而给定 $n$ 个二重极点（留数固定）的亚纯二次微分构成的向量空间维数也是 $3g-3+n$；因此「唯一性」在这两个维数上是一致的。真正深刻的是存在性（Strebel 定理），它等价于给定留数的二次微分可取为 Strebel 微分。]

\begin{example}[8.10] 本例出现在 [214, Ex. 1.2.4] 中。取 $X=\C\PP^1$，即复射影直线。亚纯二次微分 $\omega(x)=x\dd x\otimes\dd x$ 在 $X\setminus\mathrm{Crit}(\omega)$ 上确定一个实射影向量场 $V_\omega$，其中 $\mathrm{Crit}(\omega)$ 是由 $\pm\sqrt{\omega}(V_\omega)\in\R$ 所定义的 $\omega$ 的临界点集。该向量场的积分曲线满足：
\[
\int_{x_0}^{x_1}\pm\sqrt{x}\,\dd x=\pm\frac{2}{3}\left(x_1^{3/2}-x_0^{3/2}\right)\in\R\,,
\]
对所有 $x_0<x_1\in\R$。\end{example}

% [原文 p.83]
当 $x_0\ne 0$ 时，存在 $x_0$ 的邻域 $U$，使得映射 $x_1\mapsto x_1^{3/2}$ 限制在 $U$ 上（对平方根函数的任一分支选取）都是单射。于是 $U$ 中含唯一一条经过点 $x_0$ 的 $\omega$ 的轨道，其参数化为
\[
x_1(t)=\left(x_0^{3/2}+t\right)^{2/3},\qquad t\in\R\,.
\]
当 $x_0=0$ 时，$x_0$ 的邻域 $U$ 中含三条从 $0$ 出发的轨道。它们的参数化为
\[
x_{1,k}(t)=t\,\ee^{2\pi\ii k/3},\qquad t\in\R_+,\ k\in\{0,1,2\}\,.
\]
最后，在 $x=\infty$ 的邻域内，取局部坐标 $y=x^{-1}$，二次微分 $\omega$ 有五阶极点：
\[
\omega(y)=y^{-1}\dd(y^{-1})\otimes\dd(y^{-1})=y^{-5}\dd y\otimes\dd y\,.
\]
图 8.5 描述了 $x=0$ 与 $x=\infty$ 两点附近的轨道结构。

\Fig[0.74\linewidth]{figures/fig_8_5.png}{8.5}{例 8.10 中微分 $\omega$ 的轨道在 $x=0$（左）与 $x=\infty$（右）附近的形态。}

% [Tutor Note: 极点阶数可核验：由 $x=1/y$ 得 $\dd x=-\dd y/y^2$，故 $\omega=x\,\dd x^{\otimes 2}=\frac{1}{y}\cdot\frac{\dd y^2}{y^4}=y^{-5}\dd y^2$，即 $x=\infty$ 处恰为五阶极点，与正文一致。另注意 $x=0$ 是 $\omega=x\dd x^2$ 的单零点，故由 $0$ 出发恰有三条轨道，方向为 $2\pi k/3$（$k=0,1,2$）。]

\section{谱曲线}

给定一个黎曼曲面以及一族全纯微分，我们可联系一个分歧覆叠，即所谓的谱曲线。谱曲线由 Hitchin 在 [197, 198] 中引入，当时是为了对稳定向量丛的模空间作阿贝尔化。2.2 节中关于分歧覆叠的例 2.17 给出了谱曲线的局部模型。为我们的目的，采用如下版本的这一概念：

% [原文 p.84]
\begin{definition}[8.11] 设 $X$ 是亏格 $g\ge 2$ 的光滑复代数曲线，$K=\Omega^1_X$ 是 $X$ 上的典范线丛。次数为 $r$ 的谱曲线或谱覆盖是指嵌入余切丛 $T^*X$ 中的完备光滑代数曲线 $\Sigma$，使得投影 $\pi:T^*X\to X$ 是次数为 $r$ 的分歧覆叠：
\begin{equation}
\begin{array}{ccc}
f:\ \Sigma & \longrightarrow & T^*X\\[2pt]
 & & \big\downarrow{\scriptstyle\pi}\\[2pt]
 & & X\,.
\end{array}
\tag{8.1}
\end{equation}
\end{definition}

设 $\eta\in\mathrm{H}^0(T^*X,\pi^*K)$ 是 $T^*X$ 上的重言 1-形式，使得 $-\dd\eta$ 是 $T^*X$ 上的标准辛形式（$\eta$ 也称为刘维尔形式（Liouville form））。

\begin{example}[8.12] 对全纯二次微分 $q_2\in\mathrm{H}^0(X,K^{\otimes 2})$，考虑方程
\[
\eta^{\otimes 2}+q_2=0
\]
的解。解集构成谱曲线 $\Sigma\subset T^*X$，其分歧点位于 $q_2$ 的零点处。$q_2$ 的零点的阶数等于分歧指数减 1。因此若 $q_2$ 的所有零点都是单零点，则 $\Sigma$ 是简单分歧覆叠。\end{example}

% [Tutor Note: 例 8.12 的方程严格说应在 $\Sigma\subset T^*X$ 上取拉回：$\eta$ 是 $\mathrm{H}^0(T^*X,\pi^*K)$ 中的重言形式，$q_2$ 是 $X$ 上的二次微分，二者都须经 $\pi$ 拉回到 $\Sigma$ 才成为同一丛 $\pi^*K^{\otimes 2}|_\Sigma$ 的截面。另一处易错是辛形式的符号约定：当 $\eta=p\,\dd q$ 时标准辛形式为 $-\dd\eta=\dd q\wedge\dd p$，恰与原文「$-\dd\eta$ 是标准辛形式」相符。]

我们可以推广这个例子。考虑全纯微分 $k$-形式构成的元组所成的向量空间，其中 $k$ 从 1 变到 $r$：
\[
\mathcal{B}:=\bigoplus_{i=1}^{r}\mathrm{H}^0(X,K^{\otimes i})\,.
\]
利用 Riemann–Roch 公式可以证明，$\mathcal{B}$ 的维数等于 $r^2(g-1)+1$。一个点 $q=(q_1,\dots,q_r)\in\mathcal{B}$ 通过特征方程
\[
\eta^{\otimes r}+\sum_{i=1}^{r}q_i\eta^{\otimes(r-i)}=0
\]
给出一个谱曲线。另一种看法是，曲线 $\Sigma$ 可视为整体截面
\[
\eta^{\otimes r}+\sum_{i=1}^{r}\pi^*q_i\,\eta^{\otimes(r-i)}\in\mathrm{H}^0(T^*X,\pi^*K^{\otimes r})
\]
的零点集。

若谱曲线 $\Sigma$ 只有单分歧点，则称点 $q=(q_1,\dots,q_r)\in\mathcal{B}$ 是一般的。

\begin{remark}[8.13] 设 $X=\Sigma=\C$，$\pi:\Sigma\to X$，$\pi(z)=z^2$。这是一个分歧覆叠，唯一的分歧点为 $z=0$。如图 8.6 考虑 $X$ 中两条同伦的环路 $\gamma_1$ 与 $\gamma_2$。$\gamma_1$ 的两个提升都是闭的，而 $\gamma_2$ 的提升没有一个是闭的，所以它们不可能是同伦的。\end{remark}

因此，非分歧覆叠的路径提升性质（命题 2.4）对分歧覆叠不再成立。后文我们将看到，谱网络可以提供解决这一问题的途径（见 13.7 节）。

% [原文 p.85]
\Fig[0.72\linewidth]{figures/fig_8_6.png}{8.6}{分歧覆叠中不同伦的路径提升。}

\section{Strebel 微分}

在定义 8.8 中，若叶状结构 $\hat{\mathcal{F}}_\omega$ 的所有临界叶都是正则的，即 $\hat{\mathcal{F}}_\omega$ 的每片叶的一端在 $\omega$ 的零点、另一端在 $\omega$ 的极点，则称黎曼曲面上的亚纯二次微分 $\omega$ 为一般的。这一般情形使得我们可以定义黎曼曲面的理想三角剖分（见图 8.4）。现在考虑非一般的情形；特别地，我们关心 $\omega$ 是如下意义下所谓 Strebel 微分的情形：

\begin{definition}[8.14] 设 $\bar X$ 是亏格 $g\ge 2$ 的闭黎曼曲面，并固定 $\bar X$ 的一个有限标记点集 $P=\{s_1,\dots,s_n\}$，其中 $n\in\Z_{>0}$，令 $X=\bar X\setminus P$ 为相应的开黎曼曲面。又设 $\omega$ 是 $X$ 上的亚纯二次微分，在 $P$ 的各点处有二重极点，$\hat{\mathcal{F}}_\omega$ 是 $\omega$ 的无向叶状结构。若叶状结构 $\hat{\mathcal{F}}_\omega$ 没有终止于穿孔（即 $\omega$ 的极点）的叶，且所有叶都是紧的，则称 $\omega$ 为 Strebel（二次）微分。\end{definition}

由该定义可以看出，对 Strebel 微分，叶状结构 $\hat{\mathcal{F}}_\omega$ 的叶要么是闭轨道，要么是鞍连接。Strebel 微分确实存在，这是 Strebel 的一个基本定理（参见 [336, Chap. VI]）。

% [Tutor Note: 定义 8.14 的关键在于「无叶终止于穿孔」与「所有叶紧」两条：前者排除了绕穿孔螺旋的轨道，后者排除了非闭的非紧轨道。结合定理 8.9 可见，这种非一般情形恰好是 $n$ 个二重极点、留数取正实数的情形；若把留数 $a_i$ 换成 $-a_i$（如例 8.15 中的负留数），只需把 $\omega$ 换成 $-\omega$，叶状结构不变。]

\begin{example}[8.15] 本例取自 [205, Sect. 3.3]。设 $X$ 是球面去掉 3 点，考虑由下式局部描述的亚纯二次微分 $\omega$：
\[
\omega=-\frac{m_\infty^2 z^2-(m_\infty^2+m_0^2-m_1^2)z+m_0^2}{z^2(z-1)^2}\,\dd z^2\,,
\]
其中 $m_i\in\R$，$i=0,1,\infty$。注意 $\omega$ 是 Strebel 微分，在穿孔 $\{0,1,\infty\}$ 处分别有二重极点，留数分别为 $-m_0^2$、$-m_1^2$、$-m_\infty^2$。\end{example}

% [原文 p.86]
\Fig[0.40\linewidth]{figures/fig_8_7.png}{8.7}{由例 8.15 得到的两类不同的叶状结构 $\hat{\mathcal{F}}_\omega$。上方称为“分子 I”（molecule I），在 $m_\infty<m_0+m_1$ 时得到；下方称为“分子 II”（molecule II），在 $m_\infty>m_0+m_1$ 时得到。}

% ============================================================
% 第 9 章
% ============================================================

% [原文 p.87]
\chapter{高阶泰希米勒空间}

\ChapterAbstract{泰希米勒空间的代数进路提示我们更一般地考虑 $\pi_1(S_g)$ 到 $\PSL_2(\R)$ 以外的李群的表示。这自然地引导我们考虑特征簇以及高阶泰希米勒空间的概念：后者是特征簇的连通分支，且其中全由离散且忠实的表示组成。这类空间的两个主要例子是后续各章将研究的 Hitchin 分支与极大分支。}

\section{特征簇}

特征簇处于表示论、几何与数学物理的交汇处。它们研究流形基本群的表示论，以代数的方式编码若干几何对象，如克莱因几何与平坦联络。我们对这一丰富的理论作一介绍，并把注意力限于曲面。

对亏格 $g\ge 2$ 的闭、连通、可定向拓扑曲面 $S_g$，我们想研究基本群表示的集合，即群同态 $\rho:\pi_1(S_g)\to G$，其中 $G$ 是李群。事实上，更一般地可以取任意有限生成群 $\Gamma$ 的表示集 $\Hom(\Gamma,G)$，但我们主要关注基本群表示的情形（以保持与几何的联系）。如同表示论中的惯例，我们只关心表示的同构类。若一个表示是另一个表示被某个群元素共轭所得，就把它们等同。因此我们关心商
\[
\mathcal{E}(S_g,G):=\Hom(\pi_1(S_g),G)/G\,,
\]
其中 $G$ 通过共轭作用。

% [原文 p.88]
在下文讨论的适当条件下，空间 $\Hom(\pi_1(S_g),G)/G$ 含有一个具有簇结构的子空间，通常称为特征簇。与几何联系最著名的情形之一是 $G=\PSL_2(\R)$，即双曲平面的保向等距群。我们将看到，在此情形可以挑出特征簇的一个连通分支，它可等同于曲面 $S_g$ 的泰希米勒空间。

为此，我们需要如式 (4.4) 所引入的基本群的生成元与关系表示。然而对闭曲面 $S_g$，这一表示简单地是
\begin{equation}
\pi_1(S_g)=\left\langle a_1,b_1,\dots,a_g,b_g\ \middle|\ \prod_{i=1}^{g}[a_i,b_i]=1\right\rangle,
\tag{9.1}
\end{equation}
其中 $[a_i,b_i]=a_ib_ia_i^{-1}b_i^{-1}$ 是换位子。$\pi_1(S_g)$ 到连通约化实李群 $G$ 的全部表示所成之集 $\Hom(\pi_1(S_g),G)$ 可自然地等同于 $G^{2g}$ 中满足代数方程 $\prod_i[a_i,b_i]=1$ 的 $2g$-元组 $(a_1,b_1,\dots,a_g,b_g)$ 构成的子集。群 $G$ 通过共轭作用在 $\Hom(\pi_1(S_g),G)$ 上：
\[
(g\cdot\rho)=g\rho(\gamma)g^{-1}\,,
\]
其中 $g\in G$，$\rho\in\Hom(\pi_1(S_g),G)$，$\gamma\in\pi_1(S_g)$。

% [Tutor Note: 原文此式印作「$(g\cdot\rho)=g\rho(\gamma)g^{-1}$」，左端漏写自变量，正确形式应为 $(g\cdot\rho)(\gamma)=g\rho(\gamma)g^{-1}$（即先作共轭再取值）。判断依据：左端 $g\cdot\rho$ 是表示，右端是其在一个元素 $\gamma$ 处的值，二者须在同一元素上求值才相等。]

将此作用限制到约化表示的子空间 $\Hom^{\mathrm{red}}(\pi_1(S_g),G)$ 上可以保证轨道空间是 Hausdorff 的。这里，所谓约化表示是指：它与 $G$ 的李代数中的伴随表示复合后分解为不可约表示之和。当 $G$ 是代数群时，这等价于 $\pi_1(S_g)$ 在 $G$ 中的像的 Zariski 闭包是约化群。定义 $\pi_1(S_g)$ 到 $G$ 的约化表示的模空间为该轨道空间
\[
\mathcal{R}(S_g,G):=\Hom^{\mathrm{red}}(\pi_1(S_g),G)/G\,.
\]

[166] 证明了该空间具有以表示的稳定化子为指标、由实解析簇构成的分层，因此 $\mathcal{R}(S_g,G)$ 通常被称为特征簇或表示簇。

\begin{theorem}[9.1] 模空间 $\mathcal{R}(S_g,G)$ 具有实解析簇的结构；若 $G$ 是代数群则为代数簇，若 $G$ 是复群则为复簇。\end{theorem}

\begin{remark}[9.2] 对任意有限生成群 $\Gamma$ 与李群 $G$，也考虑并研究表示 $\rho:\Gamma\to G$ 的特征簇。$G$-特征簇是 Hausdorff 拓扑空间。一般地，它们不是流形，因为可能有奇点；然而它们总是局部可缩的。\end{remark}

\begin{remark}[9.3] 表示是约化的这一条件可视为 GIT 商语境下的稳定性条件。更多细节见 [330, Sect. 7]。\end{remark}

\begin{example}[9.4] 特征簇最简单的例子是交换群 $G=\U(1)=\mathbb{S}^1$ 的情形。此时特征簇可等同于曲面的雅可比簇。

% [原文 p.89]
事实上，共轭作用是平凡的，且所有表示都是约化的，故 $\mathcal{R}(S_g,\U(1))=\mathcal{E}(S_g,\U(1))=\Hom(\pi_1(S_g),\U(1))$。\end{example}

\begin{example}[9.5] 对 $G=\SU(2)$，特征簇 $\mathcal{R}(S_g,\SU(2))$ 是带奇点的辛流形。其第 $k$ 层量子化可用于定义一个拓扑量子场论（TQFT）。相应向量空间 $V_k(S_g)$ 的维数由所谓的 Verlinde 公式给出。此外，$V_k(S_g)$ 给出映射类群 $\mathrm{Mod}(S_g)$ 的有限维表示（见 [86, 255]）。\end{example}

3.4 节讨论的 Riemann–Hilbert 对应在曲面 $S_g$ 的情形可重述如下：

\begin{theorem}[9.6] 空间 $\mathcal{E}(S_g,G)$ 同构于 $S_g$ 上平坦主 $G$-丛的模空间模去规范变换。\end{theorem}

这给出了特征簇的微分几何描述。

\section{连通分支}

特征簇具有引人入胜的拓扑，因为即便是关于连通分支数目的最简单的拓扑问题，对一般李群而言也是高度非平凡的。高阶泰希米勒理论研究那些其中所有表示都离散且忠实的连通分支，从而指向这些表示的几何解释。

主拓扑不变量以如下方式给 $\mathcal{R}(S_g,G)$ 的每个连通分支联系 $\pi_1(G)$ 中的一个元素：考虑 $\rho:\pi_1(S_g)\to G$，并取 $\pi_1(S_g)$ 的一组生成元 $(a_1,b_1,\dots,a_g,b_g)$。对每个生成元，在万有覆叠 $\tilde G$ 中选取提升 $\tilde\rho(a_i)$ 与 $\tilde\rho(b_i)$。由于 $\prod_i[a_i,b_i]=1$（见式 (9.1)），我们看出
\[
e(\rho):=\prod_{i=1}^{g}[\tilde\rho(a_i),\tilde\rho(b_i)]\in\tilde G
\]
在覆叠映射 $\pi:\tilde G\to G$ 下投影到 $1\in G$。因此可将 $e(\rho)$ 等同于 $\pi_1(G)=\pi^{-1}(\{1\})$ 的一个元素。容易验证 $e(\rho)$ 不依赖于生成元的选取或提升的选取。它是联系于表示 $\rho$ 的离散不变量，因而在包含 $\rho$ 的连通分支上为常数。

% [Tutor Note: $e(\rho)$ 的良定义性可这样核对：换一个提升只在 $\tilde G$ 上乘以中心元 $c\in\ker(\pi)=\pi_1(G)$，而中心元与一切元素交换，故 $[cg_1,cg_2]=[g_1,g_2]$，即 $e(\rho)$ 与提升选取无关；此外 $\pi(e(\rho))=\prod_i[\rho(a_i),\rho(b_i)]=1$，正说明 $e(\rho)\in\pi^{-1}(\{1\})$。至于与生成元选取无关，需要用换位子恒等式逐条验证（可用 Nielsen 变换生成 $\mathrm{Aut}\,\pi_1(S_g)$ 来归约）。]

等价地，主拓扑不变量 $e(\rho)$ 是由 Riemann–Hilbert 对应（见第 3 章）给出的、与 $\rho$ 相联系的平坦丛的欧拉类。注意这个欧拉类实际上度量了把 $\rho$ 提升为 $\pi_1(S_g)$ 到万有覆叠 $\tilde G$ 中的表示的障碍。由不变量 $e(\rho)$ 取固定值 $d$ 的那些表示构成的子簇 $\mathcal{R}_d(S_g,G)\subset\mathcal{R}(S_g,G)$ 是不相交的闭子空间，但未必是连通分支。

% [原文 p.90]
所有连通分支的描述与精确计数是近年来受到大量关注的问题。

对连通紧半单李群 $G$，对每个 $d\in\pi_1(G)$，子簇 $\mathcal{R}_d(S_g,G)$ 都非空且连通。这可以通过在曲面 $S_g$ 上固定一个复结构看出，因为这样就能在所得的、底曲线为 $S_g$ 的黎曼曲面 $X$ 上引入全纯丛方法。事实上，Ramanathan [307] 把空间 $\mathcal{R}_d(S_g,G)$ 等同于 $X$ 上拓扑类为 $d$ 的多稳定全纯主 $G^{\C}$-丛的模空间 $\mathcal{M}_d(X,G)$，其中 $G^{\C}$ 表示群 $G$ 的复化，并证明了后者实际上是连通的。

对复半单李群 $G$，全纯方法仍可给出空间 $\mathcal{R}_d(S_g,G)$ 对每个 $d\in\pi_1(G)$ 的连通性，办法是把它们等同于多稳定 $G$-希格斯丛的模空间 $\mathcal{M}_d(X,G)$；所谓 $G$-希格斯丛是指由全纯主 $G$-丛与伴随丛扭以 $X$ 的典范线丛 $T^*X$ 所得丛的一个全纯截面组成的对。在复群情形下，空间 $\mathcal{R}_d(S_g,G)$ 与 $\mathcal{M}_d(X,G)$ 的等同是 Hitchin [198] 与 Donaldson [116] 对 $G=\SL_2(\C)$、Corlette [96] 与 Simpson [331] 对一般 $G$ 的合作成果。空间 $\mathcal{M}_d(X,G)$ 的连通性是用 Hitchin 在 [198, 199] 中引入的莫尔斯理论方法得到的。此外，Goldman [167]（对 $G=\SL_2(\C)$ 与 $G=\PSL_2(\C)$）以及 Li [247]（对任意半单复李群 $G$）也提供了证明 $\mathcal{M}_d(X,G)$ 连通的其他技巧。

我们把上述讨论总结为：

\begin{theorem}[9.7] 设 $G$ 是连通半单复李群或紧李群。则特征簇 $\mathcal{R}(S,G)$ 的连通分支恰由主拓扑不变量给出。\end{theorem}

注意 $\pi_1(G)\cong\pi_1(K)$，其中 $K$ 表示 $G$ 的极大紧子群；并且由 Myers 定理，$\pi_1(K)$ 是有限的。因此我们得到有限个连通分支。

对一般实李群，我们仍有有限个连通分支，但计数要困难得多。可能有多个分支具有相同的主拓扑不变量，而 $\pi_1(G)$ 中并非所有值都能取到，正如我们将在 11.1 节埃尔米特李群的情形看到的那样。

在泰希米勒情形，我们有
\[
\pi_1(\PSL_2(\R))\cong\pi_1(\SO(2))\cong\Z\,,
\]
但只有绝对值不超过 $|\chi(S_g)|=2g-2$ 的值可以取到（这来自将在 11.1 节讨论的 Milnor–Wood 不等式）。在此情形，结果是每个值恰好被取到一次，故有 $4g-3$ 个连通分支。泰希米勒空间是欧拉数为 $2g-2$（最大值）的那个连通分支。平凡表示位于不变量 $0$ 的连通分支中。Hitchin 在 [199] 中用莫尔斯理论技巧得到了一个基本的计数结果：

% [原文 p.91]
\begin{theorem}[9.8] $\PSL_n(\R)$-特征簇在 $n$ 为奇数时有 3 个连通分支，在 $n>2$ 为偶数时有 6 个连通分支。\end{theorem}

对实李群的情形，连通分支的精确计数是在非阿贝尔霍奇理论领域内实现的，那里有大量技巧可用，从而能够用多稳定 $G$-希格斯丛的模空间给出 $G$-特征簇的复解析描述。

对 $G=\GL_n(\R)$，黎曼曲面 $X$ 上的希格斯丛是指一个全纯丛 $E$ 连同全纯 1-形式 $\Phi\in\mathrm{H}^0(X,\End(E)\otimes K)$，其中 $K$ 表示 $X$ 的典范丛。非阿贝尔霍奇理论以如下方式联系一个（多稳定的）希格斯丛与一个平坦联络（即特征簇中的一个点）。给定 $E$ 上的埃尔米特度量 $h$，可考虑联络
\[
\Phi+\dd_A+\Phi^{*h}\,,
\]
其中 $\dd_A$ 是由 $h$ 确定的 $E$ 上的陈联络，$\Phi^{*h}$ 是 $\Phi$ 的埃尔米特共轭。

完全非平凡的事实是：存在 $h$ 的选取使得该联络是平坦的。表述曲率消失的方程称为 Hitchin 方程：
\[
F(\dd_A)+[\Phi\wedge\Phi^{*h}]=0\,,
\]
其中 $F(\dd_A)$ 表示陈联络的曲率。

这一美妙的理论超出了本书的范围；我们推荐 [176] 与 [343] 作为入门，并推荐原始工作 [96]、[116]、[198]、[331]。至于当李群是复单李群的实形式时这些分支计数的最近且详细的综述，我们推荐 [62]。

\section{高阶泰希米勒空间}

从 Labourie [242] 的工作出发，出现了一个基本想法：用分支中表示的动力学性质来描述特征簇的连通分支。在这一视角下，定义泰希米勒空间的如下推广：

\begin{definition}[9.9] 高阶泰希米勒空间是特征簇 $\mathcal{R}(S_g,G)$ 的若干连通分支之并，其中全由离散且忠实的表示组成。\end{definition}

「高泰希米勒空间」一词（后来演变为「高阶泰希米勒空间」）源自 Fock 与 Goncharov [129] 的工作：他们在一般分裂实半单约化李群的语境下，对 Lusztig [251] 的全正性概念发展了一种更代数几何化的进路，并定义了到这些群中的正基本表示。

% [原文 p.92]
在 [179, 180] 中，Guichard 与 Wienhard 引入了李群 $G$ 上 $\Theta$ 正结构的概念，并证明它给出 $G$-特征簇连通分支的一个第二拓扑不变量。最近，Beyrer–Guichard–Labourie–Pozzetti–Wienhard [51] 证明了主拓扑不变量与第二拓扑不变量一起构成高阶泰希米勒空间的完全不变量，从而集此课题数十年研究之大成。利用 $\Theta$ 正性这一核心概念，目前已知有四类主要的高阶泰希米勒空间：
\begin{itemize}[leftmargin=2em]
\item Hitchin 分支，对应分裂实群 $G$；
\item 极大分支，对应管型埃尔米特李群 $G$；
\item $G$-特征簇的连通分支，对应局部同构于 $\SO(p,q)$ 的群 $G$；
\item 例外型的连通分支，即群 $G$ 为 $F_4,E_6,E_7,E_8$ 型且限制根系为 $F_4$ 型的情形。
\end{itemize}

特别地，Hitchin 分支全由离散且忠实的表示组成，这一点由 Labourie 证明：他在 [242] 中引入了 Anosov 表示的概念，并使用了动力系统的技巧。Fock 与 Goncharov 在 [129] 中也建立了这一结论：他们构造了分裂实李群的全部正表示，并证明它们是忠实的、离散的且正双曲的。当群 $G$ 是管型埃尔米特李群时，关于极大分支的断言由 Burger、Iozzi、Labourie 与 Wienhard 在 [69] 与 [70] 中建立；而当李群局部同构于 $\SO(p,q)$ 时，该结论由 Beyrer 与 Pozzetti [52] 得到。

% [Tutor Note: 上述第三、四类中「局部同构于 $\SO(p,q)$」与「$F_4,E_6,E_7,E_8$ 型且限制根系为 $F_4$ 型」是两类不同的例外来源：$E_6,E_7,E_8$ 的分裂实形式之限制根系为 $F_4$ 型（因而与 $F_4$ 型同属「例外型」），而 $\SO(p,q)$ 属于经典群。原文把二者分列为第 3、4 条，翻译时照录，不要合并。]

后文我们将详细讨论高阶泰希米勒空间的两类主要族：第 10 章的 Hitchin 分支与第 11 章的极大分支。$\Theta$ 正性的一般概念则在第 12 章研究。
